\documentclass[11pt,oneside]{book}

\usepackage[T1]{fontenc}
\usepackage[utf8]{inputenc}
\usepackage{lmodern}
\usepackage{microtype}
\usepackage[a4paper,margin=28mm,headheight=14pt]{geometry}
\usepackage{amsmath,amssymb,amsthm,amscd,mathtools,bm}
\usepackage{booktabs,longtable,array,tabularx}
\usepackage{enumitem}
\usepackage{graphicx}
\usepackage{caption}
\usepackage{xcolor}
\usepackage{url}
\usepackage[unicode,colorlinks=true,linkcolor=blue!55!black,citecolor=blue!55!black,urlcolor=blue!55!black]{hyperref}
\usepackage[nameinlink,noabbrev]{cleveref}
\usepackage{fancyhdr}
\usepackage{listings}
\usepackage{etoolbox}

\IfFileExists{corpus_metadata.tex}{\input{corpus_metadata.tex}}{%
  \newcommand{\CorpusCount}{65}%
  \newcommand{\RepoCommit}{\texttt{unknown}}%
  \newcommand{\RepoCommitDate}{unknown}%
  \newcommand{\CorpusRoot}{\texttt{Analysis/FabiusFunction/docs}}%
}

\pagestyle{fancy}
\fancyhf{}
\fancyhead[L]{Closing the Lagrange--Rvachev loop}
\fancyhead[R]{\thepage}
\setlength{\parindent}{1.4em}
\setlength{\parskip}{0.15em}
\setlist[itemize]{topsep=0.35em,itemsep=0.15em}
\setlist[enumerate]{topsep=0.35em,itemsep=0.15em}
\allowdisplaybreaks

\newtheorem{theorem}{Theorem}[chapter]
\newtheorem{proposition}[theorem]{Proposition}
\newtheorem{lemma}[theorem]{Lemma}
\newtheorem{corollary}[theorem]{Corollary}
\newtheorem{conjecture}[theorem]{Conjecture}
\newtheorem{problem}[theorem]{Problem}
\theoremstyle{definition}
\newtheorem{definition}[theorem]{Definition}
\newtheorem{example}[theorem]{Example}
\newtheorem{algorithm}[theorem]{Algorithm}
\theoremstyle{remark}
\newtheorem{remark}[theorem]{Remark}

\DeclareMathOperator{\upf}{up}
\DeclareMathOperator{\supp}{supp}
\DeclareMathOperator{\ord}{ord}
\DeclareMathOperator{\rank}{rank}
\DeclareMathOperator{\ran}{ran}
\DeclareMathOperator{\tr}{tr}
\DeclareMathOperator{\diag}{diag}
\DeclareMathOperator{\spanop}{span}
\DeclareMathOperator{\dist}{dist}
\newcommand{\R}{\mathbb R}
\newcommand{\C}{\mathbb C}
\newcommand{\N}{\mathbb N}
\newcommand{\Z}{\mathbb Z}
\newcommand{\Pp}{\mathcal P}
\newcommand{\Iop}{\mathcal I}
\newcommand{\Aop}{\mathcal A}
\newcommand{\Sop}{\mathcal S}
\newcommand{\Eop}{\mathcal E}
\newcommand{\Top}{\mathcal T}
\newcommand{\M}{\mathcal M}
\newcommand{\vtwo}{\nu_2}
\newcommand{\dd}{\,\mathrm d}
\newcommand{\e}{\mathrm e}
\newcommand{\iu}{\mathrm i}
\newcommand{\norm}[1]{\left\lVert#1\right\rVert}
\newcommand{\abs}[1]{\left\lvert#1\right\rvert}
\newcommand{\set}[1]{\left\{#1\right\}}
\newcommand{\angles}[1]{\left\langle#1\right\rangle}
\newcommand{\restricted}{\mathbin{\big|}}
\newcommand{\trans}[1]{\tau_{#1}}
\newcommand{\RepoPath}[1]{\texttt{#1}}

\lstdefinestyle{code}{
  basicstyle=\ttfamily\small,
  columns=fullflexible,
  breaklines=true,
  frame=single,
  showstringspaces=false
}
\lstset{style=code}

\title{\textbf{Closing the Lagrange--Rvachev Loop}\\[0.35em]
\Large Exact polynomial synthesis by undilated dyadic translates of the Rvachev \(\operatorname{up}\)-function, interpolation factorizations, and coefficient-space projectors}
\author{Research report prepared for Vladimir Reshetnikov}
\date{29 August 2026\\Repository snapshot \RepoCommit}

% ed.: the document is book class, which provides no abstract environment, so the
% abstract below raised "Environment abstract undefined" and its text was dropped
% from the build. Supplied here in the standard report-class shape.
\newenvironment{abstract}%
  {\cleardoublepage\thispagestyle{plain}\begin{center}\bfseries Abstract\end{center}\quotation}%
  {\endquotation\clearpage}

% ed.: \Sha (the Dirac comb) is used in the weighted-Poisson appendix but never
% defined. Rendered as an upright sans-serif Cyrillic Sha built from three rules,
% which needs no Cyrillic font to be installed.
\providecommand{\Sha}{\mathord{\vcenter{\hbox{%
  \rule[-0.05ex]{0.09em}{1.15ex}\kern0.22em
  \rule[-0.05ex]{0.09em}{1.15ex}\kern0.22em
  \rule[-0.05ex]{0.09em}{1.15ex}%
  \kern-0.62em\rule[-0.05ex]{0.62em}{0.09em}}}}}

\begin{document}
\frontmatter
\maketitle

\begin{abstract}
The repository already contains an exact theorem that synthesizes a polynomial on a bounded interval from finitely many shifted copies of Rvachev's compactly supported function \(u=\operatorname{up}\).  The purpose of this report is not to restate that theorem in isolation, but to expose its transform-theoretic mechanism, specialize it to Lagrange cardinal polynomials on \([-1,1]\), and compose it with the corpus's Lagrange-interpolation expansions of \(u\).  The composition closes the apparent circle
\[
  u \longrightarrow \text{Lagrange polynomials}
  \longrightarrow \text{shifted copies of }u.
\]

With the normalization \(\supp u=[-1,1]\) and \(\int u=1\), no dilation of \(u\) is required.  If \(p\) has degree at most \(d\), put \(m=2^d\), \(h=m^{-1}\), and let
\[
  M(z)=\int_{-1}^{1}\e^{zx}u(x)\dd x
      =\prod_{r=1}^{\infty}\frac{\sinh(z/2^r)}{z/2^r}.
\]
Then on \([-1,1]\)
\[
  p(x)=h\sum_{\substack{k\in\Z\\ |k|<2m}}
  \bigl[M(D)^{-1}p\bigr](kh)\,u(x-kh).
\]
The inverse operator is finite on \(\mathcal P_d\).  Its monomial images form an Appell sequence
\[
  \sum_{n\ge0}A_n(x)\frac{z^n}{n!}=\frac{\e^{xz}}{M(z)},
\]
whose cumulants have the closed Bernoulli form
\[
  \kappa_{2r}=\frac{2^{2r-1}B_{2r}}{r(2^{2r}-1)},\qquad \kappa_{2r+1}=0.
\]
Consequently, every degree-\(d\) Lagrange cardinal \(\ell_j\) has the explicit finite representation
\[
  \ell_j(x)=\sum_{|k|<2^{d+1}} C_{kj}\,u\!\left(x-\frac{k}{2^d}\right),
  \qquad
  C_{kj}=2^{-d}\bigl[M(D)^{-1}\ell_j\bigr]\!\left(\frac{k}{2^d}\right).
\]

The Fourier explanation is sharp and purely dyadic.  At a reciprocal-lattice point,
\[
 \ord_{\omega=2\pi m\nu}\widehat u(\omega)=\vtwo(m\nu)+1.
\]
Thus the shift lattice \(m^{-1}\Z\) has exact Strang--Fix order \(\vtwo(m)+1\): its density is not what controls polynomial reproduction; only the power of two dividing \(m\) does.  The canonical choice \(m=2^d\) is the smallest positive integer that reproduces every polynomial of degree at most \(d\).

At a fixed interpolation level, substitution gives no mysterious new approximation operator.  It factors the ordinary interpolation projector exactly.  If \(X=(x_0,\ldots,x_d)\), \(R_{ik}=u(x_i-k/m)\), and \(C\) is the matrix above, then
\[
  RC=I_{d+1}.
\]
Hence \(C\) is an explicit Bernoulli--Appell right inverse of a rectangular up-collocation matrix, while \(P=CR\) is an oblique coefficient-space projector.  With the minimal interval-wide index set \(|k|<2m\), its spectrum is
\[
  1\quad(d+1\text{ times}),\qquad
  0\quad(4m-d-2\text{ times}).
\]
For data \(y_j=f(x_j)\), the synthesized loop equals the Lagrange interpolant \(\mathcal I_Xf\) identically.  Therefore its convergence, Lebesgue constant, and Runge failures are exactly those of the original interpolation scheme; the smooth atoms do not regularize an unstable projector because their coefficient cancellations reconstruct that projector exactly.

For a convergent nested or telescoping Lagrange expansion, the substitution yields a level-grouped double series of shifted copies of \(u\), and its partial sum through level \(N\) is precisely \(\mathcal I_Nu\).  The report develops this identity, proves the collocation and projector results, derives explicit dyadic inter-level null vectors, gives low-degree closed forms, reports numerical checks, and formulates several conjectures about coefficient measures, asymptotic obliquity, and stable node families.  Claims described as ``new in this report'' mean that they were not found in the recursively scanned repository corpus; no claim of priority over the entire external literature is intended.
\end{abstract}

\tableofcontents
\listoftables
\listoffigures

\chapter*{Principal conclusions}
\addcontentsline{toc}{chapter}{Principal conclusions}
The main conclusions can be read at three mutually reinforcing levels.

\paragraph{Transform level.}
The infinite sinc product has zeros of increasingly high multiplicity at increasingly even reciprocal frequencies.  For the undilated kernel, passing from integer shifts to shifts by \(2^{-d}\) raises the Strang--Fix order from one to \(d+1\).  This is why a fixed physical atom of radius one can reproduce degree-\(d\) polynomials when its centers are placed on the grid \(2^{-d}\Z\).

\paragraph{Polynomial level.}
The coefficient polynomial is not the target polynomial itself.  It is the finite deconvolution \(q=M(D)^{-1}p\).  Bernoulli numbers enter through the cumulants of \(u\); Bell polynomials package the inverse series; and the resulting monomial images \(A_n=M(D)^{-1}x^n\) form an Appell family.  For a Lagrange cardinal, one computes the cardinal once, applies this finite differential operator once, and samples the result on the dyadic grid.

\paragraph{Operator level.}
Composing Lagrange interpolation with polynomial synthesis gives an exact factorization
\[
   \Iop_X=\Sop_m\Eop_{d,m}\Iop_X.
\]
On sample vectors this says \(RC=I\).  On coefficient vectors it produces the idempotent \(CR\).  On functions it says that every truncation of the closed loop is just the original interpolant in a different coordinate system.  The equality settles convergence and Runge questions immediately and identifies the genuinely new questions as coefficient stability, null-space geometry, and limiting coefficient measures.

\mainmatter

\chapter{Corpus, scope, and notation}
\label{ch:corpus}

\section{Repository scope}
The analysis used the recursively scanned directory \CorpusRoot\ at repository commit \RepoCommit\ (commit date \RepoCommitDate).  The snapshot contains \CorpusCount\ TeX files.  The scan includes top-level notes, archived expositions, papers, and the newer frontier drafts in subdirectories.  A path-by-path manifest, extracted title, leading section headings, line count, and a reproducible lexical source map appear in \cref{app:corpus}; the machine-readable CSV and the extraction script accompany the report.

The directly relevant starting point is \RepoPath{Up\_Polynomial\_Synthesis.tex}.  The corpus already goes substantially beyond the naive statement that ``some compactly supported shifts span some polynomials'': it supplies an exact finite synthesis mechanism.  The present report therefore concentrates on the nonduplicative layer requested in the problem:
\begin{enumerate}
  \item specialize the synthesis map to arbitrary Lagrange cardinal systems on \([-1,1]\);
  \item identify the Fourier/Strang--Fix reason for the canonical dyadic denominator \(2^d\);
  \item compose the map with the Lagrange representation of \(u\), both at a fixed level and as a telescoping series;
  \item analyze the resulting collocation right inverse, coefficient projector, dyadic null vectors, convergence, and numerical stability.
\end{enumerate}

The report uses the corpus's usual normalization of the Rvachev function.  Changes of Fourier convention merely move factors of \(2\pi\); all multiplicity and polynomial-reproduction statements are invariant after the corresponding change of reciprocal lattice.

\IfFileExists{lagrange_source_map.tex}{\input{lagrange_source_map.tex}}{}

\section{The normalized Rvachev function}
Let
\[
 u(x)=\operatorname{up}(x)
\]
be the even, nonnegative, infinitely differentiable function characterized by
\begin{equation}
 \supp u=[-1,1],\qquad \int_{\R}u(x)\dd x=1,
 \label{eq:up-normalization}
\end{equation}
and the Fourier product
\begin{equation}
 \widehat u(\omega)
 :=\int_{\R}u(x)\e^{-\iu\omega x}\dd x
 =\prod_{r=1}^{\infty}\operatorname{sinc}\!\left(\frac{\omega}{2^r}\right),
 \qquad \operatorname{sinc}z:=\frac{\sin z}{z}.
 \label{eq:up-fourier-product}
\end{equation}
Equivalently, \(u\) is the density of the almost surely convergent random series
\begin{equation}
 X=\sum_{r=1}^{\infty}2^{-r}U_r,
 \qquad U_r\sim\operatorname{Unif}[-1,1]
 \quad\text{independently}.
 \label{eq:up-random-series}
\end{equation}
Its support radius is \(\sum_{r\ge1}2^{-r}=1\).  Splitting off the first summand gives the familiar integral and differential equations
\begin{align}
 u(x)&=\int_{2x-1}^{2x+1}u(t)\dd t,\label{eq:up-integral}\\
% ed.: was a single "\", which escaped the newline instead of ending the
% align row; amsmath then reported "Multiple \label's" and dropped
% eq:up-integral.
 u'(x)&=2u(2x+1)-2u(2x-1).\label{eq:up-differential}
\end{align}
In particular \(u(0)=1\), and \(u\) together with all derivatives vanishes at \(\pm1\).

\section{Polynomials, grids, and shifts}
Write \(\Pp_d\) for the real polynomials of degree at most \(d\).  For a positive integer \(m\), put
\[
 h=\frac1m,
 \qquad
 K_m:=\set{k\in\Z:\abs{k}<2m}.
 \label{eq:Km}
\]
The index set \(K_m\) is the minimal fixed lattice window whose translates can meet \([-1,1]\) nontrivially: if \(x\in[-1,1]\) and \(u(x-kh)\ne0\), then \(|k|<2m\).  Retaining \(k=\pm2m\) is harmless, but those two atoms vanish identically on the closed target interval and merely add two zero columns to every matrix.

For a coefficient vector \(c=(c_k)_{k\in K_m}\), define the restricted synthesis operator
\begin{equation}
 (\Sop_m c)(x):=\sum_{k\in K_m}c_k u(x-kh),
 \qquad -1\le x\le1.
 \label{eq:synthesis-operator}
\end{equation}
No argument of \(u\) is rescaled in \cref{eq:synthesis-operator}.  The atom is always the same physical radius-one function; only its center changes.

\section{Lagrange notation}
Let \(X=(x_0,\ldots,x_d)\) be pairwise distinct nodes in \([-1,1]\).  The cardinal polynomials and interpolation projector are
\begin{align}
 \ell_j^X(x)&:=\prod_{\substack{0\le r\le d\\r\ne j}}
 \frac{x-x_r}{x_j-x_r},\label{eq:lagrange-cardinal}\\
% ed.: was a single "\" here, which escaped the newline instead of ending the
% align row; both labels then fell in one row and eq:lagrange-cardinal was
% clobbered, leaving the two \cref's to it undefined.
 (\Iop_X f)(x)&:=\sum_{j=0}^{d}f(x_j)\ell_j^X(x).
 \label{eq:lagrange-projector}
\end{align}
When the node set is clear, the superscript is omitted.  The map from samples \(y=(y_0,\ldots,y_d)^T\) to its interpolating polynomial is denoted \(L_Xy\); thus \(L_X(f|_X)=\Iop_Xf\).

\chapter{Moment, cumulant, and Bernoulli--Appell structure}
\label{ch:appell}

\section{Moment generating function}
The moment generating function of \(u\) is
\begin{equation}
 M(z):=\int_{-1}^{1}\e^{zx}u(x)\dd x
 =\prod_{r=1}^{\infty}\frac{\sinh(z/2^r)}{z/2^r}.
 \label{eq:mgf}
\end{equation}
It is an even entire function with \(M(0)=1\).  Write
\begin{equation}
 M(z)=\sum_{n=0}^{\infty}\mu_n\frac{z^n}{n!},
 \qquad
 \log M(z)=\sum_{n=1}^{\infty}\kappa_n\frac{z^n}{n!}.
 \label{eq:moments-cumulants}
\end{equation}
Here \(\mu_n=\int x^nu(x)\dd x\), and \(\kappa_n\) is the \(n\)-th cumulant of the random variable in \cref{eq:up-random-series}.

\begin{proposition}[Bernoulli cumulants]
\label{prop:cumulants}
All odd cumulants vanish, and for \(r\ge1\),
\begin{equation}
 \kappa_{2r}=\frac{2^{2r-1}B_{2r}}{r(2^{2r}-1)}.
 \label{eq:cumulants}
\end{equation}
In particular,
\[
 \mu_0=1,\qquad \mu_2=\frac19,\qquad
 \kappa_4=-\frac{2}{225},\qquad \mu_4=\frac{19}{675}.
\]
\end{proposition}

\begin{proof}
The classical expansion
\[
 \log\frac{\sinh w}{w}
 =\sum_{r=1}^{\infty}
 \frac{2^{2r-1}B_{2r}}{r(2r)!}\,w^{2r}
\]
follows by integrating the Bernoulli expansion of \(\coth w-w^{-1}\).  Substitute \(w=z/2^j\), sum over \(j\ge1\), and use
\[
 \sum_{j=1}^{\infty}2^{-2rj}=\frac1{2^{2r}-1}.
\]
Comparison with \cref{eq:moments-cumulants} proves \cref{eq:cumulants}.  The displayed moments follow from the standard moment--cumulant relations.
\end{proof}

\section{The inverse moment operator}
Because \(M(0)=1\), the reciprocal has a formal and locally convergent expansion
\begin{equation}
 \frac1{M(z)}=\sum_{n=0}^{\infty}\lambda_n\frac{z^n}{n!}.
 \label{eq:inverse-moments}
\end{equation}
On \(\Pp_d\), define
\begin{equation}
 \Aop_d:=M(D)^{-1}
estricted_{\Pp_d}
 =\sum_{n=0}^{d}\lambda_n\frac{D^n}{n!},
 \qquad D:=\frac{\dd}{\dd x}.
 \label{eq:inverse-moment-operator}
\end{equation}
This is a finite triangular automorphism of \(\Pp_d\).  Since \(M\) is even, \(M(-D)=M(D)\), so no sign convention survives in the final formulas.

There are three complementary closed descriptions.
\begin{enumerate}
\item \emph{Cumulant/Bell form.}  From \cref{eq:cumulants},
\begin{equation}
 \frac{\e^{xz}}{M(z)}
 =\exp\!\left(xz-\sum_{r\ge1}\kappa_{2r}\frac{z^{2r}}{(2r)!}\right),
 \label{eq:appell-bell-generating}
\end{equation}
so its coefficients are complete exponential Bell polynomials evaluated at
\((x,-\kappa_2,0,-\kappa_4,0,\ldots)\).

\item \emph{Bernoulli product form.}  Since
\[
 \frac{w}{\sinh w}=\frac{2w\e^w}{\e^{2w}-1}
 =\sum_{n=0}^{\infty}B_n\!\left(\frac12\right)\frac{(2w)^n}{n!},
\]
we have, modulo \(D^{d+1}\),
\begin{equation}
 \Aop_d=
 \prod_{r=1}^{\infty}
 \left[
   \sum_{n=0}^{d}B_n\!\left(\frac12\right)
   2^{n(1-r)}\frac{D^n}{n!}
 \right]\pmod{D^{d+1}}.
 \label{eq:bernoulli-product-operator}
\end{equation}
The infinite product is computationally finite to every prescribed degree after collecting powers of \(D\).

\item \emph{Appell form.}  Define \(A_n\) by
\begin{equation}
 \sum_{n=0}^{\infty}A_n(x)\frac{z^n}{n!}
 =\frac{\e^{xz}}{M(z)}.
 \label{eq:appell-definition}
\end{equation}
Then
\begin{equation}
 A_n=\Aop_n(x^n),\qquad M(D)A_n=x^n,
 \qquad A_n'=nA_{n-1},
 \qquad A_n(-x)=(-1)^nA_n(x).
 \label{eq:appell-properties}
\end{equation}
\end{enumerate}

\begin{proposition}[Appell recurrence]
\label{prop:appell-recurrence}
With \(A_0=1\),
\begin{equation}
 A_{n+1}(x)=xA_n(x)-
 \sum_{r=1}^{n}\binom nr\kappa_{r+1}A_{n-r}(x).
 \label{eq:appell-recurrence}
\end{equation}
Only odd values of \(r\) contribute, because \(\kappa_{r+1}=0\) when \(r+1\) is odd.
\end{proposition}

\begin{proof}
Differentiate \cref{eq:appell-definition} with respect to \(z\).  The logarithmic derivative of its right-hand side is
\[
 x-\frac{M'(z)}{M(z)}
 =x-\sum_{r\ge0}\kappa_{r+1}\frac{z^r}{r!}.
\]
Equating exponential-generating-function coefficients gives \cref{eq:appell-recurrence}.
\end{proof}

The first terms, generated exactly by the accompanying Python code, are as follows.
\IfFileExists{appell_table.tex}{\input{appell_table.tex}}{
\begin{align*}
A_0(x)&=1,&A_1(x)&=x,&A_2(x)&=x^2-\frac19,\\
A_3(x)&=x^3-\frac{x}{3},&
A_4(x)&=x^4-\frac23x^2+\frac{31}{675}.
\end{align*}}

\section{Deconvolution of a general polynomial}
If
\[
 p(x)=\sum_{n=0}^{d}a_nx^n,
\]
then
\begin{equation}
 q(x):=(\Aop_dp)(x)=\sum_{n=0}^{d}a_nA_n(x).
 \label{eq:q-appell-linear-combination}
\end{equation}
Equivalently, using \cref{eq:inverse-moment-operator},
\begin{equation}
 q(x)=\sum_{r=0}^{d}\lambda_r\frac{p^{(r)}(x)}{r!}.
 \label{eq:q-derivative-form}
\end{equation}
This is the ``one finite Bernoulli--Appell deconvolution'' referred to throughout the report.  It is independent of the synthesis grid.  The grid enters only in the subsequent sampling of \(q\).

\chapter{Exact dyadic spectral multiplicities}
\label{ch:spectral}

\section{Zeros of the infinite sinc product}
For a nonzero integer \(n\), let \(\vtwo(n)\) be the exponent of two in \(n\).  The key observation is that several factors of \cref{eq:up-fourier-product} vanish at a sufficiently even reciprocal frequency.

\begin{theorem}[Exact dyadic zero multiplicity]
\label{thm:zero-multiplicity}
For positive integers \(m\) and nonzero integers \(\nu\),
\begin{equation}
 \ord_{\omega=2\pi m\nu}\widehat u(\omega)
 =\vtwo(m\nu)+1
 =\vtwo(m)+\vtwo(\nu)+1.
 \label{eq:zero-multiplicity}
\end{equation}
\end{theorem}

\begin{proof}
Write \(m\nu=2^q a\) with \(a\) odd.  At \(\omega_0=2\pi m\nu\), the \(r\)-th factor in \cref{eq:up-fourier-product} is
\[
 \operatorname{sinc}\!\left(\frac{\omega_0}{2^r}\right)
 =\operatorname{sinc}\!\left(2^{q-r+1}\pi a\right).
\]
It has a simple zero exactly for \(1\le r\le q+1\).  For \(r\ge q+2\), the argument is a nonintegral dyadic multiple of \(\pi\), hence the factor is nonzero.  The product of the first \(q+1\) simple zeros with a nonzero analytic tail has order \(q+1\).
\end{proof}

\begin{corollary}[Sharp Strang--Fix order]
\label{cor:strang-fix-order}
For the shift lattice \(h\Z=m^{-1}\Z\), the minimum zero order of \(\widehat u\) over nonzero reciprocal-lattice points \(2\pi m\nu\) is
\begin{equation}
 s(m):=\vtwo(m)+1.
 \label{eq:strang-fix-order}
\end{equation}
Consequently the undilated translates \(u(\,\boldsymbol\cdot-k/m)\) have exact polynomial-reproduction degree \(\vtwo(m)\).  In particular, \(m=2^d\) gives order \(d+1\) and reproduces \(\Pp_d\); no smaller positive integer does so for all degree-\(d\) polynomials.
\end{corollary}

\begin{remark}[Density is not the invariant]
A grid with a huge odd denominator still has only first-order reciprocal zeros and therefore only constant reproduction.  Replacing \(m\) by \(2m\) raises the order by exactly one.  Thus the relevant arithmetic parameter is \(\vtwo(m)\), not \(m\) itself.  This 2-adic formulation is both sharper and conceptually cleaner than saying merely that ``finer dyadic grids reproduce higher degrees.''
\end{remark}

\section{A weighted Poisson identity}
The bridge from zero multiplicity to explicit coefficients is a polynomially weighted Poisson formula.

\begin{theorem}[Polynomial comb equals continuous convolution]
\label{thm:poisson-polynomial}
Let \(m\in\N\), \(h=m^{-1}\), and \(d\le\vtwo(m)\).  For every \(q\in\Pp_d\),
\begin{equation}
 h\sum_{k\in\Z}q(kh)u(x-kh)
 =\int_{\R}q(t)u(x-t)\dd t
 =\bigl[M(-D)q\bigr](x).
 \label{eq:weighted-poisson}
\end{equation}
For each fixed \(x\), the sum is finite because \(u\) is compactly supported.
\end{theorem}

\begin{proof}
It is enough to prove the identity for \(q(t)=t^r\), \(0\le r\le d\).  Introduce a formal parameter \(z\) and apply Poisson summation to
\(g_z(t)=\e^{zt}u(x-t)\).  Its Fourier transform in \(t\) is
\[
 \widehat g_z(\omega)
 =\e^{(z-\iu\omega)x}M(-z+\iu\omega).
\]
Hence, formally near \(z=0\),
\begin{equation}
 h\sum_{k\in\Z}\e^{zkh}u(x-kh)
 =\sum_{\nu\in\Z}
 \e^{(z-2\pi\iu m\nu)x}
 M(-z+2\pi\iu m\nu).
 \label{eq:poisson-egf}
\end{equation}
For \(\nu\ne0\), \cref{thm:zero-multiplicity} says that
\(M(-z+2\pi\iu m\nu)\) vanishes to order at least \(d+1\) at \(z=0\).  Therefore these terms make no contribution to the coefficients of \(z^0,\ldots,z^d\).  The \(\nu=0\) term is \(\e^{zx}M(-z)\), the exponential generating function of
\(\int t^r u(x-t)\dd t\).  Comparing coefficients proves the first equality.  The second follows from
\[
 \int q(t)u(x-t)\dd t
 =\int q(x-y)u(y)\dd y
 =M(-D)q(x).
\]
A distributional proof, obtained by differentiating the Dirac comb in frequency, gives the same result without the formal parameter.
\end{proof}

\begin{corollary}[Sharp failure one degree higher]
\label{cor:sharp-failure}
If \(d=\vtwo(m)\), then \cref{eq:weighted-poisson} cannot hold for every \(q\in\Pp_{d+1}\).  At odd reciprocal harmonics the first nonvanishing alias derivative occurs exactly at order \(d+1\).
\end{corollary}

\begin{proof}
At \(\nu=1\), \cref{thm:zero-multiplicity} gives zero order exactly \(d+1\), so the coefficient of \(z^{d+1}\) in the nonzero-harmonic part of \cref{eq:poisson-egf} is not identically zero.  A suitable degree-\(d+1\) monomial detects it.
\end{proof}

\chapter{Finite polynomial synthesis on \texorpdfstring{$[-1,1]$}{[-1,1]}}
\label{ch:synthesis}

\section{The synthesis theorem in transform form}
We now combine continuous deconvolution with the exact quadrature identity.

\begin{theorem}[Undilated dyadic synthesis]
\label{thm:polynomial-synthesis}
Let \(p\in\Pp_d\), choose any positive integer \(m\) with \(\vtwo(m)\ge d\), put \(h=m^{-1}\), and define
\begin{equation}
 q:=\Aop_dp=M(D)^{-1}p.
 \label{eq:q-definition}
\end{equation}
Then, for every \(x\in[-1,1]\),
\begin{equation}
 p(x)=h\sum_{k\in K_m}q(kh)u(x-kh).
 \label{eq:polynomial-synthesis}
\end{equation}
The canonical minimal denominator is \(m=2^d\).  Formula \cref{eq:polynomial-synthesis} uses literal translates of the same function \(u\); it contains no dilation.
\end{theorem}

\begin{proof}
By definition, \(M(D)q=p\).  Since \(u\) is even, \(M(-D)=M(D)\).  Applying \cref{thm:poisson-polynomial} gives
\[
 h\sum_{k\in\Z}q(kh)u(x-kh)
 =M(-D)q(x)=p(x).
\]
For \(x\in[-1,1]\), a term can be nonzero only if \(|x-kh|<1\), which implies \(|k|<2m\).  Thus the infinite sum reduces to \(K_m\).
\end{proof}

\begin{remark}[Local rather than global equality]
A nonconstant polynomial cannot equal a compactly supported finite sum on all of \(\R\).  The equality is asserted on the prescribed physical interval \([-1,1]\).  More generally, on any compact target interval \([a,b]\), the same infinite-lattice identity truncates to the finitely many centers within distance one of \([a,b]\).
\end{remark}

\begin{remark}[Why the factor \(h\) is necessary]
At degree zero, \cref{eq:polynomial-synthesis} becomes
\[
 h\sum_{k\in\Z}u(x-kh)=1.
\]
Thus \(\sum_k u(x-kh)=m\), as required by Poisson summation and \(\int u=1\).  Omitting \(h\) would already fail for constants.
\end{remark}

\section{A canonical polynomial analysis map}
Define
\begin{equation}
 \Eop_{d,m}:\Pp_d\longrightarrow\R^{K_m},
 \qquad
 (\Eop_{d,m}p)_k
 :=h(\Aop_dp)(kh).
 \label{eq:analysis-map}
\end{equation}
Then \cref{thm:polynomial-synthesis} is the left-inverse relation
\begin{equation}
 \Sop_m\Eop_{d,m}=J_d,
 \label{eq:synthesis-analysis-left-inverse}
\end{equation}
where \(J_d:\Pp_d\hookrightarrow C[-1,1]\) is inclusion.  In particular, \(\Eop_{d,m}\) is injective.

The map is canonical in a useful, limited sense.  Many coefficient vectors synthesize the same polynomial on a bounded interval.  The choice \cref{eq:analysis-map} is distinguished by three properties:
\begin{enumerate}
\item it is obtained from the translation-invariant convolution inverse \(M(D)^{-1}\);
\item it samples one polynomial \(q\) on the whole lattice with the natural quadrature weight \(h\);
\item it is compatible with the moment and Appell structure of \(u\), without solving a node-dependent linear system.
\end{enumerate}
It is not, in general, the Euclidean minimum-norm coefficient vector.

\section{Explicit monomial identities}
Taking \(p=x^n\) in \cref{thm:polynomial-synthesis} gives
\begin{equation}
 x^n=h\sum_{k\in K_m}A_n(kh)u(x-kh),
 \qquad \vtwo(m)\ge n,
 \quad -1\le x\le1.
 \label{eq:monomial-synthesis}
\end{equation}
The first cases are
\begin{align}
 1&=h\sum_{k\in K_m}u(x-kh),\label{eq:constant-synthesis}\\
 x&=h\sum_{k\in K_m}(kh)u(x-kh),\label{eq:linear-synthesis}\\
 x^2&=h\sum_{k\in K_m}\left((kh)^2-\frac19\right)u(x-kh),\label{eq:quadratic-synthesis}\\
 x^3&=h\sum_{k\in K_m}\left((kh)^3-\frac{kh}{3}\right)u(x-kh),\label{eq:cubic-synthesis}\\
 x^4&=h\sum_{k\in K_m}\left((kh)^4-\frac23(kh)^2+\frac{31}{675}\right)u(x-kh).
 \label{eq:quartic-synthesis}
\end{align}
These identities display directly why raw polynomial samples are insufficient: the moment corrections \(-1/9\), \(-x/3\), and so forth undo the smoothing by \(u\).

\section{Affine target intervals}
For a polynomial identity on \([a,b]\), no rescaling of the atom is needed.  One merely keeps all centers \(kh\) whose radius-one supports intersect the interval:
\begin{equation}
 p(x)=h\sum_{\substack{k\in\Z\\ a-1<kh<b+1}}
 (\Aop_dp)(kh)u(x-kh),\qquad a\le x\le b.
 \label{eq:arbitrary-interval}
\end{equation}
If instead one elects to scale the physical kernel to radius \(\rho\), then the moment operator and reciprocal lattice scale covariantly.  The present problem has \(\rho=1\), for which \cref{eq:polynomial-synthesis} already gives shifts of the original \(\operatorname{up}\)-function exactly as requested.

\chapter{Lagrange cardinals as finite sums of shifted \texorpdfstring{$u$}{up}}
\label{ch:lagrange-synthesis}

\section{General node set}
Apply \cref{thm:polynomial-synthesis} to each cardinal \(\ell_j^X\in\Pp_d\).

\begin{theorem}[Finite synthesis of every Lagrange cardinal]
\label{thm:lagrange-synthesis}
Let \(X=(x_0,\ldots,x_d)\subset[-1,1]\) have distinct nodes, let \(m\) satisfy \(\vtwo(m)\ge d\), and put \(h=m^{-1}\).  Define
\begin{equation}
 q_j^X:=\Aop_d\ell_j^X,
 \qquad
 C_{kj}^{X,m}:=h q_j^X(kh),
 \quad k\in K_m.
 \label{eq:C-coefficients}
\end{equation}
Then
\begin{equation}
 \boxed{
 \ell_j^X(x)=\sum_{k\in K_m}C_{kj}^{X,m}u(x-kh),
 \qquad -1\le x\le1.}
 \label{eq:lagrange-synthesis-boxed}
\end{equation}
For the canonical choice \(m=2^d\),
\begin{equation}
 C_{kj}^{X,2^d}
 =2^{-d}\left[M(D)^{-1}\ell_j^X\right]\!\left(\frac{k}{2^d}\right),
 \qquad |k|<2^{d+1}.
 \label{eq:canonical-cardinal-coefficients}
\end{equation}
\end{theorem}

\begin{proof}
This is \cref{thm:polynomial-synthesis} with \(p=\ell_j^X\).
\end{proof}

If \(\ell_j^X(x)=\sum_{r=0}^{d}v_{rj}x^r\), then the Appell form is
\begin{equation}
 q_j^X(x)=\sum_{r=0}^{d}v_{rj}A_r(x),
 \qquad
 C_{kj}^{X,m}=h\sum_{r=0}^{d}v_{rj}A_r(kh).
 \label{eq:C-appell}
\end{equation}
This makes the computational separation clean: invert the Vandermonde matrix once to obtain \(v_{rj}\), tabulate \(A_r(kh)\) once, and multiply.

\section{Vandermonde--Appell matrix formula}
Let
\begin{equation}
 V_X=(x_j^r)_{0\le j,r\le d},
 \qquad
 B_{m,d}=\bigl(hA_r(kh)\bigr)_{k\in K_m,\ 0\le r\le d}.
 \label{eq:V-B-matrices}
\end{equation}
With columns indexed by cardinals, their monomial coefficient matrix is \(V_X^{-1}\), under the convention that a sample vector \(y\) gives monomial coefficients \(V_X^{-1}y\).  Therefore
\begin{equation}
 C^{X,m}=B_{m,d}V_X^{-1}.
 \label{eq:C-matrix-factorization}
\end{equation}
Equation \cref{eq:C-matrix-factorization} is the finite Bernoulli--Appell transform of the inverse Vandermonde matrix.  It avoids evaluating products in \cref{eq:lagrange-cardinal} separately when many cardinals are required.

\section{Equispaced and dyadic combs}
For the degree-\(d\) equispaced nodes
\[
 x_j=-1+\frac{2j}{d},\qquad 0\le j\le d,
\]
the cardinal has the factorial form
\begin{equation}
 \ell_j(x)=
 \frac{(-1)^{d-j}}{j!(d-j)!}
 \prod_{\substack{0\le r\le d\\r\ne j}}
 \left(\frac d2(x+1)-r\right).
 \label{eq:equispaced-cardinal}
\end{equation}
For the dyadic comb
\begin{equation}
 X_N:=\set{\frac{j}{N}:-N\le j\le N},
 \qquad N=2^n,
 \label{eq:dyadic-comb}
\end{equation}
the degree is \(d=2N\), and
\begin{equation}
 \ell_{N,j}(x)=
 \frac{(-1)^{N-j}}{(N+j)!(N-j)!}
 \prod_{\substack{-N\le r\le N\\r\ne j}}(Nx-r).
 \label{eq:dyadic-cardinal}
\end{equation}
The canonical undilated-up synthesis grid for this cardinal degree has
\begin{equation}
 m=2^d=2^{2N}.
 \label{eq:double-exponential-m}
\end{equation}
This large denominator is a theorem-friendly minimal choice, not a recommendation for direct high-degree numerical computation.  It reflects the exact zero multiplicity needed to reproduce every degree-\(2N\) polynomial with the unscaled kernel.

\section{Low-degree identities}
\begin{example}[Two endpoint cardinals]
For nodes \((-1,1)\), \(d=1\), \(m=2\), and
\[
 \ell_-(x)=\frac{1-x}{2},\qquad \ell_+(x)=\frac{1+x}{2}.
\]
Because \(\Aop_1\) fixes linear polynomials,
\begin{align}
 \ell_-(x)&=\sum_{|k|<4}\frac{2-k}{8}\,u\!\left(x-\frac k2\right),\label{eq:linear-cardinal-minus}\\
 \ell_+(x)&=\sum_{|k|<4}\frac{2+k}{8}\,u\!\left(x-\frac k2\right).
 \label{eq:linear-cardinal-plus}
\end{align}
\end{example}

\begin{example}[Cardinals at \(-1,0,1\)]
For \(d=2\), \(m=4\),
\[
 \ell_{-1}=\frac{x^2-x}{2},\qquad
 \ell_0=1-x^2,\qquad
 \ell_1=\frac{x^2+x}{2}.
\]
Using \(A_2=x^2-1/9\), one obtains, for \(|k|<8\),
\begin{align}
 C_{k,-1}&=\frac{k^2}{128}-\frac{k}{32}-\frac1{72},\\
 C_{k,0}&=\frac5{18}-\frac{k^2}{64},\label{eq:quadratic-center-coeff}\\
 C_{k,1}&=\frac{k^2}{128}+\frac{k}{32}-\frac1{72}.
\end{align}
In particular,
\begin{equation}
 1-x^2=\sum_{|k|<8}
 \left(\frac5{18}-\frac{k^2}{64}\right)
 u\!\left(x-\frac k4\right).
 \label{eq:quadratic-up-identity}
\end{equation}
\end{example}

\begin{example}[First nontrivial dyadic-comb interpolant of \(u\)]
At the five nodes \(-1,-\tfrac12,0,\tfrac12,1\), the values are
\[
 0,\ \frac12,\ 1,\ \frac12,\ 0.
\]
The even quartic interpolant is
\begin{equation}
 p_4(x)=1-\frac73x^2+\frac43x^4.
 \label{eq:p4}
\end{equation}
Its deconvolution is
\begin{equation}
 q_4(x)=\Aop_4p_4(x)
 =\frac43x^4-\frac{29}{9}x^2+\frac{2674}{2025}.
 \label{eq:q4}
\end{equation}
With \(m=16\),
\begin{equation}
 p_4(x)=\frac1{16}\sum_{|k|<32}
 q_4\!\left(\frac{k}{16}\right)
 u\!\left(x-\frac{k}{16}\right).
 \label{eq:p4-self-representation}
\end{equation}
This is already a closed-loop identity: the polynomial on the left is the Lagrange interpolant of \(u\), while every atom on the right is a shifted copy of \(u\).
\end{example}

\chapter{Closing the loop at a fixed interpolation level}
\label{ch:fixed-loop}

\section{Direct substitution}
Let \(f\) be any function defined on the nodes \(X=(x_0,\ldots,x_d)\).  Substitute \cref{eq:lagrange-synthesis-boxed} into \cref{eq:lagrange-projector}:
\begin{align}
 (\Iop_Xf)(x)
 &=\sum_{j=0}^{d}f(x_j)\sum_{k\in K_m}C_{kj}^{X,m}u(x-kh)\notag\\
 &=\sum_{k\in K_m}b_k^{X,m}(f)u(x-kh),
 \label{eq:fixed-loop-substitution}
\end{align}
where
\begin{equation}
 b_k^{X,m}(f)
 :=\sum_{j=0}^{d}C_{kj}^{X,m}f(x_j)
 =h\left[\Aop_d(\Iop_Xf)\right](kh).
 \label{eq:fixed-loop-coefficients}
\end{equation}
For \(f=u\), this is the requested finite self-representation of the \emph{Lagrange truncation} of \(u\):
\begin{equation}
 (\Iop_Xu)(x)
 =\sum_{k\in K_m}b_k^{X,m}(u)u(x-kh).
 \label{eq:up-interpolant-self-representation}
\end{equation}
The distinction between \(u\) and \(\Iop_Xu\) is essential.  At a finite level, the loop closes on the interpolant.  It closes on \(u\) itself only after taking a convergent limit of interpolation levels.

\begin{theorem}[Exact factorization of the interpolation projector]
\label{thm:interpolation-factorization}
For every node set \(X\) of cardinality \(d+1\) and every \(m\) with \(\vtwo(m)\ge d\),
\begin{equation}
 \boxed{\Iop_X=\Sop_m\Eop_{d,m}\Iop_X.}
 \label{eq:interpolation-factorization}
\end{equation}
Consequently, the closed-loop operator
\begin{equation}
 \Top_{X,m}:=\Sop_m\Eop_{d,m}\Iop_X
 \label{eq:closed-loop-operator}
\end{equation}
is exactly \(\Iop_X\), not merely asymptotic to it.
\end{theorem}

\begin{proof}
The range of \(\Iop_X\) is \(\Pp_d\).  The identity follows by applying \cref{eq:synthesis-analysis-left-inverse} to every polynomial in that range.
\end{proof}

\section{A commutative diagram}
The fixed-level construction is summarized by
\[
\begin{CD}
 C[-1,1] @>{\Iop_X}>> \Pp_d @>{\Eop_{d,m}}>> \R^{K_m}\\
 @V{\Iop_X}VV @V{J_d}VV @VV{\Sop_m}V\\
 \Pp_d @= \Pp_d @>{J_d}>> C[-1,1].
\end{CD}
\]
To avoid a dependency on a diagram package, the same content can be read as the pair of identities
\[
 \Sop_m\Eop_{d,m}=J_d,
 \qquad
 \Sop_m\Eop_{d,m}\Iop_X=\Iop_X.
\]
The first identity belongs to polynomial synthesis; the second is the closed loop.

\section{Immediate consequences}
Because the operators are identical, every function-space property of the loop truncation is inherited from Lagrange interpolation.

\begin{corollary}[Same error, same convergence set]
\label{cor:same-error}
For every \(f\), every \(x\in[-1,1]\), and every admissible \(m\),
\begin{equation}
 f(x)-(\Top_{X,m}f)(x)=f(x)-(\Iop_Xf)(x).
 \label{eq:same-error}
\end{equation}
For a sequence of node sets \(X_n\), the loop truncations converge pointwise, uniformly, in \(L^p\), or in any other norm if and only if the corresponding interpolation projectors do so on that function.
\end{corollary}

\begin{corollary}[Same Lebesgue function]
\label{cor:same-lebesgue}
The pointwise operator norm of the loop is the ordinary Lebesgue function
\begin{equation}
 \Lambda_X(x)=\sum_{j=0}^{d}\abs{\ell_j^X(x)},
 \label{eq:lebesgue-function}
\end{equation}
and
\begin{equation}
 \norm{\Top_{X,m}}_{C\to C}
 =\norm{\Iop_X}_{C\to C}
 =\Lambda_X:=\max_{x\in[-1,1]}\Lambda_X(x).
 \label{eq:lebesgue-constant}
\end{equation}
\end{corollary}

\begin{remark}[Smooth atoms do not regularize]
Each atom \(u(x-kh)\) is \(C^\infty\), compactly supported, and visually benign.  Nevertheless, an unstable equispaced interpolation projector remains unstable after factorization.  The large alternating coefficients reconstruct the same high-degree polynomial and therefore the same endpoint oscillation.  Any suppression of Runge's phenomenon must change the polynomial projector---for example the nodes, the interpolation conditions, or the use of least squares---rather than merely changing coordinates inside its range.
\end{remark}

\chapter{Collocation right inverses and coefficient-space projectors}
\label{ch:matrices}

\section{The rectangular collocation matrix}
Let \(r_m(x)\in\R^{K_m}\) be the atom vector
\begin{equation}
 r_m(x):=\bigl(u(x-kh)\bigr)_{k\in K_m}.
 \label{eq:atom-vector}
\end{equation}
Define the \((d+1)\times(4m-1)\) collocation matrix
\begin{equation}
 R^{X,m}_{ik}:=u(x_i-kh),
 \qquad 0\le i\le d,
 \quad k\in K_m.
 \label{eq:collocation-matrix}
\end{equation}
The cardinal synthesis matrix \(C^{X,m}\) from \cref{eq:C-coefficients} has size \((4m-1)\times(d+1)\).

\begin{theorem}[Explicit right inverse]
\label{thm:right-inverse}
For every admissible \(m\),
\begin{equation}
 \boxed{R^{X,m}C^{X,m}=I_{d+1}.}
 \label{eq:right-inverse}
\end{equation}
Thus \(R^{X,m}\) has full row rank and \(C^{X,m}\) is an explicit right inverse constructed without solving the collocation system.
\end{theorem}

\begin{proof}
The \((i,j)\)-entry of \(RC\) is
\[
 \sum_{k\in K_m}u(x_i-kh)C_{kj}^{X,m}
 =\ell_j^X(x_i)=\delta_{ij}
\]
by \cref{eq:lagrange-synthesis-boxed}.
\end{proof}

In vector form, the cardinal row is
\begin{equation}
 \bigl(\ell_0^X(x),\ldots,\ell_d^X(x)\bigr)
 =r_m(x)^TC^{X,m}.
 \label{eq:cardinal-row-factorization}
\end{equation}
For a sample vector \(y\),
\begin{equation}
 (L_Xy)(x)=r_m(x)^TC^{X,m}y.
 \label{eq:sample-factorization}
\end{equation}
This is the finite matrix form of the loop.

\section{The oblique projector}
Define
\begin{equation}
 P^{X,m}:=C^{X,m}R^{X,m}
 \quad\text{on }\R^{K_m}.
 \label{eq:coefficient-projector}
\end{equation}

\begin{theorem}[Coefficient-space projector]
\label{thm:coefficient-projector}
The matrix \(P^{X,m}\) is idempotent and has
\begin{align}
 (P^{X,m})^2&=P^{X,m},\label{eq:P-idempotent}\\
 \ran P^{X,m}&=\ran C^{X,m},\label{eq:P-range}\\
 \ker P^{X,m}&=\ker R^{X,m},\label{eq:P-kernel}\\
 \rank P^{X,m}&=d+1.\label{eq:P-rank}
\end{align}
With \(|K_m|=4m-1\), its algebraic and geometric eigenvalue multiplicities are
\begin{equation}
 \begin{array}{c|c}
 \text{eigenvalue}&\text{multiplicity}\\ \hline
 1&d+1\\
 0&4m-d-2.
 \end{array}
 \label{eq:P-spectrum}
\end{equation}
Equivalently,
\begin{equation}
 \tr P^{X,m}=d+1,
 \qquad
 \chi_P(t)=t^{4m-d-2}(t-1)^{d+1}.
 \label{eq:P-characteristic}
\end{equation}
\end{theorem}

\begin{proof}
Using \cref{eq:right-inverse},
\[
 P^2=CRCR=C(RC)R=CR=P.
\]
The range and kernel statements follow from the general algebra of a right inverse.  Since \(C\) is injective and has \(d+1\) columns, \(\rank P=d+1\).  An idempotent is diagonalizable with eigenvalues zero and one, yielding the multiplicities and characteristic polynomial.
\end{proof}

\begin{remark}[Endpoint-inclusive convention]
If one uses \(-2m\le k\le2m\), the two extreme columns of \(R\) are zero on \([-1,1]\).  The projector then gains two additional zero eigenvalues.  The rank and the nonzero spectral subspace are unchanged.
\end{remark}

\begin{remark}[Obliquity]
The projector is generally not orthogonal in the Euclidean coefficient norm.  Orthogonality would require \(C=R^\dagger\), the Moore--Penrose right inverse.  The Bernoulli--Appell inverse is selected instead by polynomial deconvolution and lattice sampling.  Its norm
\(\norm{P^{X,m}}\) measures the angle between \(\ran C\) and \(\ker R\), and is a natural stability diagnostic not visible in the function-level equality \(\Top_{X,m}=\Iop_X\).
\end{remark}

\section{All right inverses and the canonical one}
Every right inverse of \(R\) has the form
\begin{equation}
 C+N,
 \qquad RN=0.
 \label{eq:all-right-inverses}
\end{equation}
Thus node interpolation alone does not determine coefficients between nodes.  The synthesis theorem imposes the much stronger interval identity \cref{eq:lagrange-synthesis-boxed}; among interval-exact right inverses, there can still be local null combinations of translates.  The map \(C=B_{m,d}V_X^{-1}\) is canonical relative to the continuous convolution inverse \(M(D)^{-1}\), but not uniquely forced by finite-dimensional linear algebra.

\section{Sample-space and coefficient-space geometry}
The maps form the split exact sequence
\begin{equation}
 0\longrightarrow\ker R
 \longrightarrow\R^{K_m}
 \mathop{\longrightarrow}^{R}\R^{d+1}
 \longrightarrow0,
 \label{eq:split-sequence}
\end{equation}
with splitting \(C\).  Hence
\begin{equation}
 \R^{K_m}=\ran C\oplus\ker R.
 \label{eq:direct-sum}
\end{equation}
The projector \(P=CR\) keeps the canonical polynomial coefficient component and removes the collocation-null component.  A coefficient vector \(c\) and its projected vector \(Pc\) agree at the nodes after synthesis:
\begin{equation}
 R(Pc)=Rc.
 \label{eq:projected-same-data}
\end{equation}
Only when \(c\in\ran C\) is it already the Bernoulli--Appell representation of the unique interpolating polynomial.

\chapter{The Lagrange series and the level-grouped self-representation}
\label{ch:series-loop}

\section{Telescoping form of the existing expansion}
Let \(X_n\subset[-1,1]\) be a sequence of node sets, with \(|X_n|=d_n+1\), and write \(\Iop_n:=\Iop_{X_n}\).  Whenever
\begin{equation}
 \Iop_nu\longrightarrow u
 \label{eq:interpolation-convergence}
\end{equation}
in a chosen topology, one has the exact telescoping Lagrange series
\begin{equation}
 u=\Iop_0u+\sum_{n=1}^{\infty}\Delta_nu,
 \qquad
 \Delta_n:=\Iop_n-\Iop_{n-1}.
 \label{eq:lagrange-telescoping-series}
\end{equation}
Every term is a polynomial and, after expansion in the relevant cardinal bases, a finite linear combination of Lagrange polynomials.  This operator form covers the dyadic-comb and other node families considered in the corpus and cleanly separates the algebraic loop from the analytic question of whether \cref{eq:interpolation-convergence} holds.

\section{Synthesize each level separately}
For each \(n\), choose \(m_n\) with
\begin{equation}
 \vtwo(m_n)\ge \max(d_n,d_{n-1}),
 \label{eq:mn-condition}
\end{equation}
where \(d_{-1}:=-1\).  The canonical choice is
\(m_n=2^{\max(d_n,d_{n-1})}\).  Put
\begin{equation}
 a_{n,k}:=m_n^{-1}
 \left[\Aop_{\max(d_n,d_{n-1})}(\Delta_nu)\right]
 \left(\frac{k}{m_n}\right),
 \qquad |k|<2m_n,
 \label{eq:level-coefficients}
\end{equation}
with \(\Delta_0u:=\Iop_0u\).  Then
\begin{equation}
 \Delta_nu(x)=\sum_{|k|<2m_n}a_{n,k}
 u\!\left(x-\frac{k}{m_n}\right).
 \label{eq:level-synthesis}
\end{equation}
Substitution into \cref{eq:lagrange-telescoping-series} gives the closed loop
\begin{equation}
 \boxed{
 u(x)=\sum_{n=0}^{\infty}\sum_{|k|<2m_n}
 a_{n,k}u\!\left(x-\frac{k}{m_n}\right),}
 \label{eq:double-self-representation}
\end{equation}
where the series is grouped by levels.

\begin{theorem}[Partial sums are exactly the original interpolants]
\label{thm:partial-sums}
For every \(N\),
\begin{equation}
 \sum_{n=0}^{N}\sum_{|k|<2m_n}
 a_{n,k}u\!\left(x-\frac{k}{m_n}\right)
 =\Iop_Nu(x).
 \label{eq:partial-sum-equals-interpolant}
\end{equation}
Therefore \cref{eq:double-self-representation} converges in a topology exactly when \(\Iop_Nu\to u\) in that topology.
\end{theorem}

\begin{proof}
Each inner sum equals \(\Delta_nu\) by \cref{eq:level-synthesis}.  Summing through \(N\) telescopes.
\end{proof}

\begin{remark}[Order of summation matters]
The theorem licenses level-grouped summation.  It does not by itself permit flattening the pairs \((n,k)\), rearranging terms, or asserting absolute convergence of
\(\sum_{n,k}|a_{n,k}|\).  Such statements require independent coefficient estimates.  This distinction is important because interpolation can converge through large cancellations even when the total variation of its synthesized coefficient measures grows rapidly.
\end{remark}

\section{Synthesize a fixed partial sum once}
There is a second, often simpler representation.  Choose
\begin{equation}
 M_N=2^{d_N}
 \label{eq:MN}
\end{equation}
and synthesize the single polynomial \(\Iop_Nu\):
\begin{equation}
 \Iop_Nu(x)=\sum_{|k|<2M_N}b_{N,k}
 u\!\left(x-\frac{k}{M_N}\right),
 \label{eq:fixed-level-self-representation}
\end{equation}
where
\begin{equation}
 b_{N,k}=M_N^{-1}
 \left[\Aop_{d_N}(\Iop_Nu)\right]\!\left(\frac{k}{M_N}\right).
 \label{eq:fixed-level-b}
\end{equation}
The level-by-level and single-level partial sums are identical as functions but generally have different coefficient vectors and even live on different grids.  Their differences are explicit local null combinations of translates, studied next.

\section{Direct cardinal substitution at level \(N\)}
If \(X_N=(x_{N,0},\ldots,x_{N,d_N})\), then
\begin{equation}
 b_{N,k}=\sum_{j=0}^{d_N}u(x_{N,j})C_{kj}^{X_N,M_N}.
 \label{eq:b-cardinal-sum}
\end{equation}
This is precisely the operation requested in the problem statement: take the existing Lagrange representation of \(u\), replace each cardinal polynomial by its finite shifted-up representation, interchange two finite sums, and collect coefficients.  The operator analysis shows what the substitution means: it is the factorization of \(\Iop_N\) through \(\Pp_{d_N}\), the Bernoulli--Appell sampling map, and the up-synthesis map.

\chapter{Dyadic inter-level relations and local null vectors}
\label{ch:dyadic-null}

\section{Two exact representations of the same polynomial}
Fix \(p\in\Pp_d\), let \(m\) satisfy \(\vtwo(m)\ge d\), and put
\[
 q=\Aop_dp.
\]
At grid levels \(m\) and \(2m\), \cref{thm:polynomial-synthesis} gives
\begin{align}
 p(x)&=\frac1m\sum_{|k|<2m}q\!\left(\frac{k}{m}\right)
 u\!\left(x-\frac{k}{m}\right),\label{eq:coarse-rep}\\
 p(x)&=\frac1{2m}\sum_{|r|<4m}q\!\left(\frac{r}{2m}\right)
 u\!\left(x-\frac{r}{2m}\right).
 \label{eq:fine-rep}
\end{align}
Embed the coarse vector in the fine grid by placing it at even indices.  Subtracting gives a family of exact local dependencies among undilated translates.

\begin{theorem}[Dyadic quadrature-detail null vector]
\label{thm:dyadic-null}
Define \(\delta^{q,m}=(\delta_r)_{|r|<4m}\) by
\begin{equation}
 \delta_r=
 \begin{cases}
 \dfrac1{2m}q(r/2m),&r\text{ odd},\\[0.7em]
 -\dfrac1{2m}q(r/2m),&r\text{ even}.
 \end{cases}
 \label{eq:dyadic-null-coefficients}
\end{equation}
Then
\begin{equation}
 \sum_{|r|<4m}\delta_r
 u\!\left(x-\frac{r}{2m}\right)=0,
 \qquad -1\le x\le1.
 \label{eq:dyadic-null-identity}
\end{equation}
As \(q\) ranges over \(\Pp_d\), these null vectors form a \((d+1)\)-dimensional subspace.
\end{theorem}

\begin{proof}
At an even index \(r=2k\), the fine coefficient is
\((2m)^{-1}q(k/m)\), while the embedded coarse coefficient is
\(m^{-1}q(k/m)\); their difference is
\(-(2m)^{-1}q(k/m)\).  At odd indices there is no coarse coefficient, so the difference is the fine coefficient.  Subtract \cref{eq:coarse-rep} from \cref{eq:fine-rep}.  Injectivity of polynomial sampling at more than \(d\) points shows that the map \(q\mapsto\delta^{q,m}\) is injective.
\end{proof}

\begin{remark}[An undilated analogue of a detail mask]
The alternating parity in \cref{eq:dyadic-null-coefficients} resembles a wavelet detail mask, but the atoms are not dilated between levels.  The cancellation compares two exact quadrature formulas for the same convolution.  Iterating the construction produces a nested hierarchy of explicit null spaces indexed by dyadic refinement.
\end{remark}

\section{Moment cancellation of the detail measure}
Associate the signed measure
\begin{equation}
 \eta_{q,m}:=\sum_{|r|<4m}\delta_r\,\delta_{r/(2m)}.
 \label{eq:detail-measure}
\end{equation}
For \(x\in[-1,1]\), \cref{eq:dyadic-null-identity} says
\begin{equation}
 (u*\eta_{q,m})(x)=0.
 \label{eq:convolution-null-measure}
\end{equation}
When \(q\) is a monomial or an Appell polynomial, the coefficients give explicit finite identities.  These measures are not globally in the kernel of convolution by \(u\); their convolution vanishes on the target interval.  The distinction between global translate independence and local interval dependencies is therefore crucial.

\section{Only the odd aliases govern the new order}
Passing from \(m\) to \(2m\) increases every reciprocal zero multiplicity by one, but the minimum is attained at odd harmonics.  Even harmonics already carry extra multiplicity through \(\vtwo(\nu)\).  This yields a dyadic spectral filtration:
\begin{equation}
 \ord_{2\pi m\nu}\widehat u
 =\vtwo(m)+1+\vtwo(\nu).
 \label{eq:spectral-filtration}
\end{equation}
The leading alias obstruction to one more reproduced degree lies entirely in the odd reciprocal harmonics.  Higher-even harmonics enter only at correspondingly higher derivative orders.  This graded structure is invisible in a scalar Strang--Fix order and is potentially useful for refined error expansions.

\chapter{Convergence, Runge behavior, and numerical conditioning}
\label{ch:convergence}

\section{Exact-arithmetic convergence}
Let \(X_n\) be any sequence of node sets.  From \cref{thm:partial-sums}, the function-level question is completely settled:
\begin{equation}
 \norm{u-\text{loop}_N}
 =\norm{u-\Iop_nu}
 \label{eq:loop-error}
\end{equation}
for every norm in which both sides are interpreted.  Thus:
\begin{itemize}
\item if the Lagrange expansion converges uniformly, so does the closed loop;
\item if it converges only pointwise or on compact subintervals, the loop has exactly that behavior;
\item if a Runge-type endpoint blow-up occurs, it reappears without attenuation;
\item changing the admissible synthesis denominator \(m\) changes coefficients, not the exact polynomial or its error.
\end{itemize}

\section{Node geometry remains decisive}
For equispaced degree-\(d\) interpolation, the Lebesgue constant grows exponentially in \(d\); for Chebyshev--Lobatto nodes it grows logarithmically.  The compact support and infinite smoothness of \(u\) do not alter these operator norms.  Since the zero extension of \(u\) is \(C^\infty\) but not analytic at the support boundary, Chebyshev approximation is expected to be superalgebraic rather than geometrically convergent.  Equispaced interpolation has no comparable general stability guarantee, and the corpus's Runge investigations should be read directly through \cref{eq:interpolation-factorization}.

The practical recommendation is therefore simple: choose the polynomial approximation mechanism first.  Once a stable polynomial \(p_d\) has been selected---by Chebyshev interpolation, Hermite interpolation, constrained least squares, or another projector---the synthesis theorem provides an exact shifted-up coordinate representation of that \(p_d\).

\section{Internal coefficient amplification}
Although \(\Top_{X,m}=\Iop_X\) exactly, a floating-point implementation performs the factorization
\begin{equation}
 f|_X\xmapsto{\ C\ }b\xmapsto{\ r_m(x)^T\ }\Iop_Xf(x).
 \label{eq:numerical-factorization}
\end{equation}
Its forward error can be much larger than that of a stable barycentric evaluation because
\begin{equation}
 \abs{\operatorname{fl}(r^Tb)-r^Tb}
 \lesssim \gamma_{|K_m|}\sum_{k\in K_m}\abs{b_k}\abs{u(x-kh)},
 \label{eq:dot-product-error}
\end{equation}
where \(\gamma_N\approx Nu_{\rm mach}\) in the usual model.  The exact result may rely on cancellation among many terms.  Three quantities should be distinguished:
\begin{enumerate}
\item the function-space norm \(\norm{\Iop_X}=\Lambda_X\);
\item the analysis norm \(\norm{C}\), which converts samples to atom coefficients;
\item the projector obliquity \(\norm{CR}\), which measures sensitivity inside coefficient space.
\end{enumerate}
Only the first is fixed by ordinary interpolation theory.

\section{Oversampling is not by itself the instability}
For a fixed polynomial \(p\) and its deconvolution \(q\), define the coefficient measure
\begin{equation}
 \mu_{p,m}:=h\sum_{k\in K_m}q(kh)\delta_{kh}.
 \label{eq:coefficient-measure}
\end{equation}
As \(m\to\infty\), ordinary Riemann-sum estimates give
\begin{align}
 \norm{\mu_{p,m}}_{\mathrm{TV}}
 &=h\sum_{k\in K_m}\abs{q(kh)}
 \longrightarrow\int_{-2}^{2}\abs{q(t)}\dd t,
 \label{eq:l1-limit}\\
 \norm{(h q(kh))_k}_{\ell^2}
 &\sim h^{1/2}\norm{q}_{L^2(-2,2)}.
 \label{eq:l2-limit}
\end{align}
Thus the exponential number \(4m-1\) of atoms at the canonical choice \(m=2^d\) does not automatically make the \(\ell^1\) mass exponential for a \emph{fixed} polynomial.  The serious growth with \(d\) comes from the changing polynomial \(p_d\), the interpolation geometry, and the increasingly high-order inverse operator \(\Aop_d\).

\section{Stable implementation strategies}
For low and moderate degree, the following workflow is preferable to solving \(RC=I\) directly:
\begin{enumerate}
\item represent the interpolant in a numerically stable polynomial basis, preferably Chebyshev rather than monomials;
\item apply \(\Aop_d=\sum\lambda_rD^r/r!\) using stable basis differentiation and high precision when needed;
\item sample \(q\) on the dyadic grid and multiply by \(h\);
\item evaluate the atom sum in compensated or pairwise order, exploiting the evenness of \(u\) when the data are symmetric;
\item verify against direct barycentric evaluation, which computes the same polynomial by a different factorization.
\end{enumerate}
At high degree, explicitly materializing \(m=2^d\) is primarily of theoretical value.  One should instead seek compressed coefficient descriptions, multilevel cancellation formulas, or a different kernel/dilation convention when computation is the goal.

\chapter{Numerical experiments}
\label{ch:numerics}

\section{Reproducible method}
The accompanying script \RepoPath{numerical\_experiments.py} uses exact SymPy arithmetic for cumulants, inverse moments, Appell polynomials, Lagrange cardinals, and synthesis coefficients.  It tabulates \(u\) independently by inverse Fourier series on a period-four box:
\begin{equation}
 u(x)\approx\frac1{4}\sum_{|n|<N/2}
 \widehat u\!\left(\frac{\pi n}{2}\right)
 \e^{\iu\pi nx/2}.
 \label{eq:fft-tabulation}
\end{equation}
Because \(u\) is supported in \([-1,1]\), its period-four copies do not overlap.  The product \cref{eq:up-fourier-product} is truncated only after its arguments are far into the small-argument regime.  The default grid contains \(2^{18}\) points and forty-five sinc factors.

The tests are deliberately redundant:
\begin{enumerate}
\item synthesize every equispaced cardinal for degrees zero through six and compare on a dense grid;
\item form \(R\) and \(C\), then check \(RC-I\) and \((CR)^2-CR\);
\item verify the explicit quadratic and quartic closed-loop identities;
\item compare equispaced and Chebyshev--Lobatto interpolation errors for \(u\), emphasizing that these are automatically the errors of the corresponding loop truncations.
\end{enumerate}

\IfFileExists{experiment_tables.tex}{\input{experiment_tables.tex}}{
\paragraph{Build note.} Numerical tables were not generated in this fallback build; run \texttt{python numerical\_experiments.py} before compiling.}

\begin{figure}[htbp]
 \centering
 \IfFileExists{interpolation_error.pdf}{\includegraphics[width=0.82\textwidth]{interpolation_error.pdf}}{\fbox{Run the numerical script to generate the interpolation-error plot.}}
 \caption{Maximum interpolation error for two node families.  The plotted curves are simultaneously the exact-arithmetic errors of the corresponding shifted-up loop truncations, because the loop equals the interpolation projector.}
 \label{fig:interpolation-error}
\end{figure}

\section{Interpretation}
The synthesis and matrix residuals remain near the Fourier-tabulation and floating-point floor at low degrees, confirming three independent symbolic claims: the Appell correction, the exact right inverse, and projector idempotence.  Residual growth with degree is expected because the canonical grid grows as \(2^d\) and the evaluation involves increasingly severe cancellation.  It should not be interpreted as a failure of the exact theorem.

The interpolation experiment separates approximation theory from synthesis.  Chebyshev--Lobatto nodes track the smooth target stably, whereas equispaced errors eventually reflect the rapidly growing Lebesgue constant and endpoint sensitivity.  Replacing each cardinal by a smooth sum of \(u\)-atoms changes neither curve in exact arithmetic.

\chapter{Generalizations that follow without new analytic input}
\label{ch:generalizations}

\section{Any polynomial projector factors}
Lagrange interpolation is not essential to the synthesis step.  Let
\begin{equation}
 Q:C[-1,1]\longrightarrow\Pp_d
 \label{eq:any-projector}
\end{equation}
be any linear polynomial-valued operator.  Then
\begin{equation}
 Q=\Sop_m\Eop_{d,m}Q,
 \qquad \vtwo(m)\ge d.
 \label{eq:any-projector-factorization}
\end{equation}
If \(Q\) is a projector, the loop inherits its projector identity and norm.  This covers:
\begin{itemize}
\item Hermite interpolation with endpoint derivative constraints;
\item Birkhoff or lacunary interpolation whenever it is poised;
\item least-squares and discrete orthogonal projections;
\item constrained polynomial approximants designed to suppress endpoint oscillation;
\item Taylor, Chebyshev, Legendre, or other polynomial truncations.
\end{itemize}
Thus the corpus's proposal to fix higher derivatives of an interpolant can be transported into shifted-up coordinates with no additional existence theorem.  The analytic benefit, if any, comes from the modified polynomial projector; the synthesis is exact bookkeeping.

\section{Vector- and matrix-valued polynomials}
All formulas act componentwise.  If \(p:[-1,1]\to\R^s\) or \(p:[-1,1]\to\R^{r\times s}\) has polynomial entries of degree at most \(d\), then
\begin{equation}
 p(x)=h\sum_{k\in K_m}(\Aop_dp)(kh)u(x-kh),
 \label{eq:matrix-polynomial-synthesis}
\end{equation}
where \(\Aop_d\) differentiates each entry.  This may be useful for polynomial matrix approximants, transfer matrices, or simultaneous synthesis of a cardinal basis.

\section{Several variables by tensor products}
Let
\begin{equation}
 U_n(x_1,\ldots,x_n)=\prod_{j=1}^{n}u(x_j).
 \label{eq:tensor-up}
\end{equation}
On a tensor dyadic grid with denominators \(m_j\), the Fourier zeros and deconvolution operator tensorize.  A polynomial with coordinate degree at most \(d_j\) in \(x_j\) is reproduced if \(\vtwo(m_j)\ge d_j\).  The coefficient polynomial is
\begin{equation}
 \prod_{j=1}^{n}M(D_j)^{-1}p.
 \label{eq:multivariate-deconvolution}
\end{equation}
For total-degree spaces, this criterion is sufficient but not minimal; exploiting the anisotropic zero multiplicities becomes an interesting lattice-design problem.

\section{Other infinite-convolution kernels}
Suppose a compactly supported density has transform
\begin{equation}
 \widehat\phi(\omega)=\prod_{r\ge1}\widehat\rho(a_r\omega)
 \label{eq:general-infinite-convolution}
\end{equation}
with arithmetically nested zeros.  The same program applies:
\begin{enumerate}
\item count zero multiplicities on a reciprocal lattice;
\item infer the exact Strang--Fix order;
\item compute the moment inverse \(M_\phi(D)^{-1}\);
\item sample the deconvolved polynomial and truncate by support.
\end{enumerate}
The exceptional feature of \(u\) is that the powers-of-two scales make the multiplicity exactly 2-adic and the inverse cumulants explicitly Bernoulli.

\chapter{Results not found in the scanned corpus}
\label{ch:new-results}

This chapter isolates the deductions that appear to go beyond the existing repository treatment.  They are consequences of the synthesis theorem plus the transform/operator viewpoint developed here.  ``Not found'' refers to the \CorpusCount-file snapshot listed in \cref{app:corpus}; it is not a claim about all published literature.

\section{Proved deductions}
\subsection{Exact 2-adic multiplicity, not merely increasing order}
\Cref{thm:zero-multiplicity} gives the pointwise formula
\[
 \ord_{2\pi m\nu}\widehat u=\vtwo(m)+\vtwo(\nu)+1,
\]
including the extra multiplicity at even harmonics.  The scalar Strang--Fix order is only the minimum of this richer spectrum.

\subsection{Minimal denominator for undilated degree reproduction}
The canonical \(m=2^d\) is not just a convenient sufficient choice.  Among positive integer denominators, it is minimal because the exact order is \(\vtwo(m)+1\).  A denser grid with smaller 2-adic valuation cannot replace it.

\subsection{Explicit right inverse of the up-collocation matrix}
The formula
\[
 C=B_{m,d}V_X^{-1}
\]
constructs a right inverse \(RC=I\) for arbitrary distinct nodes.  It combines an inverse Vandermonde transform with universal Appell samples and requires no inversion of the generally much wider collocation matrix.

\subsection{Coefficient projector and its complete spectrum}
The matrix \(P=CR\) is an explicit oblique projector with rank \(d+1\), kernel \(\ker R\), and only eigenvalues zero and one with the multiplicities in \cref{eq:P-spectrum}.  This is the natural coefficient-space shadow of the ordinary interpolation projector.

\subsection{Dyadic detail null vectors}
\Cref{thm:dyadic-null} gives an explicit \((d+1)\)-dimensional family of local dependencies between level \(m\) and level \(2m\).  These identities explain why many coefficient descriptions can represent the same polynomial on the target interval and open a route to multilevel compression.

\subsection{Exact inheritance of Runge behavior}
The equality \(\Top_{X,m}=\Iop_X\) proves, without asymptotic estimates, that the loop has exactly the same Lebesgue function, convergence domain, and Runge failures as the original Lagrange projector.  The synthesis neither improves nor worsens exact-arithmetic approximation; it changes coordinates and may worsen floating-point cancellation.

\section{A distributional view of the loop}
For the fixed-level interpolant \(p_d=\Iop_{X_d}u\), define
\begin{equation}
 q_d=\Aop_dp_d,
 \qquad
 \mu_d=2^{-d}\sum_{|k|<2^{d+1}}
 q_d\!\left(\frac{k}{2^d}\right)\delta_{k/2^d}.
 \label{eq:mu-d}
\end{equation}
Then
\begin{equation}
 u*\mu_d=p_d\quad\text{on }[-1,1].
 \label{eq:u-mu-p}
\end{equation}
The trivial exact self-representation of \(u\) is \(u*\delta_0=u\).  The closed loop asks whether the increasingly oscillatory discrete measures \(\mu_d\) approach \(\delta_0\) in a meaningful weak topology when \(p_d\to u\).  This is a more informative limiting question than pointwise convergence of coefficients, because their lattices and dimensions change with \(d\).

\chapter{Conjectures and open problems}
\label{ch:conjectures}

\begin{conjecture}[Weak collapse to the identity atom]
\label{conj:delta-limit}
Let \(X_d\) be Chebyshev--Lobatto nodes, let \(p_d=\Iop_{X_d}u\), and define \(\mu_d\) by \cref{eq:mu-d}.  Then
\begin{equation}
 \mu_d\longrightarrow\delta_0
 \label{eq:delta-limit}
\end{equation}
in the space of compactly supported distributions.  Equivalently, for every test function \(\varphi\in C_c^\infty((-2,2))\),
\[
 2^{-d}\sum_{|k|<2^{d+1}}
 q_d(k/2^d)\varphi(k/2^d)\longrightarrow\varphi(0).
\]
The total variations \(\norm{\mu_d}_{\mathrm{TV}}\) need not remain bounded.
\end{conjecture}

\paragraph{Rationale.}
Formally, \(q_d=M(D)^{-1}p_d\) and \(M(D)^{-1}u=\delta_0\), because convolution by \(u\) multiplies Fourier transforms by \(\widehat u\).  The challenge is to justify passage from polynomial approximation to an infinite-order inverse operator whose symbol has poles at the zeros of \(\widehat u\).

\begin{conjecture}[Stability dichotomy by node family]
\label{conj:stability-dichotomy}
For Chebyshev--Lobatto interpolation of \(u\), the canonical coefficient measures \(\mu_d\) have subexponential total-variation growth:
\begin{equation}
 \limsup_{d\to\infty}\norm{\mu_d}_{\mathrm{TV}}^{1/d}=1.
 \label{eq:subexponential-tv}
\end{equation}
For equispaced interpolation, there exists a subsequence on which the corresponding total variation grows exponentially.
\end{conjecture}

\paragraph{Rationale.}
The function-level projector norms have logarithmic versus exponential behavior for these node families.  The inverse moment operator can amplify high derivatives, so the coefficient measure may be less stable than the interpolant itself; nevertheless the node dichotomy should remain visible.

\begin{conjecture}[Asymptotic weighted minimum energy]
\label{conj:min-energy}
Although \(C^{X,m}\) is generally not the Euclidean Moore--Penrose right inverse of \(R^{X,m}\), it is asymptotically minimum-norm for a translation-invariant quadrature weight as \(m\to\infty\) with \(d\) fixed.  More precisely, there should exist positive diagonal weights \(W_m\), converging to an \(L^2(-2,2)\) density, for which
\begin{equation}
 C^{X,m}=W_m^{-1}(R^{X,m})^T
 \bigl(R^{X,m}W_m^{-1}(R^{X,m})^T\bigr)^{-1}+o(1)
 \label{eq:weighted-pseudoinverse}
\end{equation}
in the operator norm on fixed-dimensional sample space.
\end{conjecture}

\paragraph{Rationale.}
The canonical coefficients are Riemann samples of the continuous deconvolution polynomial.  This suggests a variational characterization in a continuum coefficient space, but the correct energy may be Sobolev or convolution-induced rather than plain \(L^2\).

\begin{conjecture}[Generation of the polynomial local kernel]
\label{conj:null-generation}
Fix \(d\) and a sufficiently fine dyadic lattice.  Among coefficient vectors whose amplitudes are samples of polynomials of degree at most \(d\) on residue classes, the local kernel of synthesis on \([-1,1]\) is generated by iterated dyadic detail vectors of \cref{thm:dyadic-null} and boundary-invisible columns.
\end{conjecture}

\paragraph{Rationale.}
The exact zero multiplicities provide one new null layer whenever the denominator doubles.  A dimension count in a properly defined graded coefficient space should make the statement precise and possibly provable.

\begin{problem}[Refined alias-error expansion]
\label{prob:alias-expansion}
Develop an approximation formula beyond \(\Pp_d\) that retains the full multiplicity
\(\vtwo(m)+\vtwo(\nu)+1\) at each reciprocal harmonic.  The expected result is an error expansion graded simultaneously by derivative order and the 2-adic valuation of the alias index, with odd aliases supplying the leading term.
\end{problem}

\begin{problem}[Compressed exact evaluation]
\label{prob:compressed-evaluation}
The canonical degree-\(d\) formula uses \(4\cdot2^d-1\) atoms.  Find a compressed representation or multilevel algorithm that evaluates the same polynomial exactly on \([-1,1]\) without materializing all coefficients, using the dyadic null identities to cancel or aggregate levels symbolically.
\end{problem}

\begin{problem}[Hermite-stabilized loop]
\label{prob:hermite-loop}
For the derivative-constrained interpolation schemes studied elsewhere in the corpus, determine node/derivative patterns that minimize the function-level projector norm and the coefficient obliquity simultaneously.  The first objective suppresses Runge behavior; the second controls cancellation in the shifted-up realization.
\end{problem}

\chapter{Promising research directions}
\label{ch:future}

\section{A 2-adic approximation theory for infinite sinc products}
The exact multiplicity formula suggests treating reciprocal frequencies as a 2-adically filtered set.  Classical Strang--Fix theory takes the minimum zero order and discards the rest.  Retaining the valuation of each harmonic may sharpen saturation theorems, error constants, and superconvergence results for dyadic data.

\section{Coefficient measures and deconvolution topology}
The conjectural limit \(\mu_d\to\delta_0\) should be attacked in Fourier space.  For a test function \(\varphi\), one must control
\[
 \angles{\mu_d-\delta_0,\varphi}
\]
while \(q_d\) is obtained by a degree-dependent truncation of \(1/M(D)\).  Gevrey norms are natural because \(u\) is \(C^\infty\) but nowhere analytic in the relevant sense, and the infinite sinc product has a highly structured zero set.

\section{Canonical duals and frames on a bounded interval}
The columns of \(C\) behave like a dual family to the sampled translates in the finite data space.  This invites a frame-theoretic formulation: identify Hilbert spaces and weights in which the Bernoulli--Appell right inverse becomes a canonical dual, and quantify how the dual changes with the target interval and node set.

\section{Null-space bases and multiresolution compression}
The explicit parity detail vector should be iterated and orthogonalized to produce a sparse basis of local coefficient null spaces.  Such a basis could separate:
\begin{itemize}
\item the \((d+1)\)-dimensional polynomial component;
\item dyadic quadrature-detail components;
\item boundary-localized components;
\item any residual nonpolynomial local dependencies.
\end{itemize}
This decomposition is a plausible route to compressed exact synthesis and numerically stable cancellation.

\section{Alternative node systems}
Because the loop exactly inherits its projector, the strongest immediate improvements come from node selection.  The most promising comparisons are:
\begin{enumerate}
\item Chebyshev--Lobatto and Gauss--Lobatto nodes;
\item nested Leja points, which retain a telescoping structure without equispaced Runge growth;
\item dyadic nodes selected adaptively rather than taking the entire comb;
\item Hermite systems that encode the flatness of \(u\) at \(\pm1\);
\item constrained least-squares polynomial approximants with endpoint derivative conditions.
\end{enumerate}
For each family one should report both the usual Lebesgue/projection norm and the new coefficient metrics \(\norm{C}\), \(\norm{CR}\), and \(\norm{\mu_d}_{\mathrm{TV}}\).

\section{Base-\(b\) and nonuniform analogues}
Replace the binary random series by other nested uniform convolutions or compactly supported refinable densities.  The central questions are which reciprocal lattices acquire high multiplicity, which arithmetic valuation replaces \(\vtwo\), and whether the inverse cumulants retain a recognizable Bernoulli or \(q\)-series form.  This direction connects naturally to the repository's broader work on \(q\)-binomial coefficients and infinite products.

\section{Computer algebra and formal verification}
The main identities are unusually amenable to formalization:
\begin{itemize}
\item the cumulant formula is a formal power-series calculation;
\item \(M(D)^{-1}\) is finite on \(\Pp_d\);
\item the zero-multiplicity count reduces to elementary 2-adic arithmetic;
\item the matrix identities reduce to cardinal evaluation.
\end{itemize}
A Lean development could define truncated formal series, Appell polynomials, and the finite coefficient matrix exactly.  The analytic Poisson step could first be isolated as a theorem under explicit hypotheses, while all downstream polynomial and projector results would then be purely algebraic.

\chapter{Conclusion}
The requested representation exists in a particularly rigid form.  A degree-\(d\) Lagrange cardinal on \([-1,1]\) is synthesized by finitely many literal shifts of the same radius-one Rvachev function.  The coefficient calculation is
\[
 \ell_j
 \xmapsto{\ M(D)^{-1}\ }
 q_j
 \xmapsto{\text{sample at }k/2^d\text{ and multiply by }2^{-d}}
 (C_{kj})_k.
\]
The binary denominator is dictated exactly by the multiplicities of the infinite sinc product at reciprocal-lattice points.  Bernoulli cumulants, Bell polynomials, and the inverse Appell sequence describe the deconvolution explicitly.

Substitution into the Lagrange expansion of \(u\) closes the loop, but the operator interpretation prevents overclaiming: at every finite level the result is exactly the original interpolation polynomial.  This observation is powerful rather than disappointing.  It proves convergence equivalence, explains Runge failures, identifies a rectangular collocation right inverse, produces a complete coefficient-projector spectrum, and reveals dyadic null relations between levels.  The frontier is not whether the loop reproduces the interpolant---that is exact---but how its coefficient geometry behaves, whether the discrete deconvolution measures converge to \(\delta_0\), and how dyadic null spaces can be organized into a stable compressed transform.

\appendix

\IfFileExists{corpus_index.tex}{\input{corpus_index.tex}}{
\chapter{Corpus audit and source map}
\label{app:corpus}% ed.: the fallback branch carried no label, so the two \cref's
% to app:corpus were undefined whenever corpus_index.tex is absent -- which is
% always here, since this report arrived as a bare .tex with none of the
% companion files (corpus index, CSV, extraction script) its text refers to.
The generated corpus index was unavailable in this build.}

\chapter{Proof details for the weighted Poisson identity}
\label{app:poisson}

\section{Distributional derivation}
Let
\begin{equation}
 \Sha_h:=h\sum_{k\in\Z}\delta_{kh}.
 \label{eq:dirac-comb}
\end{equation}
With the Fourier convention of this report,
\begin{equation}
 \widehat{\Sha_h}=2\pi\sum_{\nu\in\Z}\delta_{2\pi\nu/h}
 \label{eq:comb-fourier}
\end{equation}
up to the standard placement of the \(2\pi\) normalization.  Multiplication of \(\Sha_h\) by \(t^r\) differentiates its Fourier transform \(r\) times.  Convolution with \(u\) multiplies these derivative deltas by \(\widehat u\).  If \(\widehat u\) has a zero of order at least \(r+1\) at every nonzero reciprocal point, all nonzero delta derivatives disappear.  The remaining origin term is exactly the Fourier transform of continuous convolution by the polynomial \(t^r\).  Linearity yields \cref{eq:weighted-poisson}.

This proof makes sharpness transparent: when one reciprocal zero has order exactly \(r\), the product with a derivative delta of order \(r\) leaves a nonzero multiple of an ordinary delta.

\section{Moment-polynomial form}
Define the discrete periodized moments
\begin{equation}
 P_r^{(m)}(x):=h\sum_{k\in\Z}(kh)^r u(x-kh).
 \label{eq:periodized-moments}
\end{equation}
For \(r\le\vtwo(m)\), \cref{thm:poisson-polynomial} gives
\begin{equation}
 P_r^{(m)}(x)=\int t^ru(x-t)\dd t
 =\sum_{j=0}^{r}\binom rj x^{r-j}(-1)^j\mu_j.
 \label{eq:periodized-moment-polynomial}
\end{equation}
Since odd moments vanish, the first examples are
\begin{align*}
 P_0^{(m)}(x)&=1,\\
 P_1^{(m)}(x)&=x,\\
 P_2^{(m)}(x)&=x^2+\frac19,\\
 P_3^{(m)}(x)&=x^3+\frac{x}{3},\\
 P_4^{(m)}(x)&=x^4+\frac23x^2+\frac{19}{675}.
\end{align*}
The Appell corrections invert this triangular relation.

\chapter{Coefficient construction in exact arithmetic}
\label{app:algorithm}

\begin{algorithm}[Cardinal-to-up transform]
\label{alg:cardinal-to-up}
Given distinct nodes \(x_0,\ldots,x_d\):
\begin{enumerate}
\item Set \(m=2^d\), \(h=1/m\), and \(K_m=\set{k\in\Z:|k|<2m}\).
\item Compute the inverse-moment coefficients through degree \(d\):
\[
 \sum_{n=0}^{d}\lambda_n\frac{z^n}{n!}
 \equiv
 \exp\!\left(-\sum_{r=1}^{\lfloor d/2\rfloor}
 \frac{2^{2r-1}B_{2r}}{r(2^{2r}-1)}
 \frac{z^{2r}}{(2r)!}\right)
 \pmod{z^{d+1}}.
\]
\item Form each cardinal \(\ell_j\), or invert the Vandermonde matrix once.
\item Apply
\[
 q_j(x)=\sum_{n=0}^{d}\lambda_n\frac{\ell_j^{(n)}(x)}{n!}.
\]
\item Output
\[
 C_{kj}=h q_j(kh),\qquad k\in K_m.
\]
\end{enumerate}
Then \(\ell_j(x)=\sum_kC_{kj}u(x-kh)\) for every \(x\in[-1,1]\).
\end{algorithm}

\section{Matrix version}
Construct the universal Appell sample matrix
\[
 B_{m,d}[k,r]=hA_r(kh).
\]
For any node set \(X\), compute its Vandermonde matrix \(V_X\), then
\[
 C^{X,m}=B_{m,d}V_X^{-1}.
\]
For data \(y\), the closed-loop coefficient vector is simply
\[
 b=C^{X,m}y.
\]
The identity \(RC=I\) is both a theorem and an exact-arithmetic regression test.

\chapter{Reproducibility files}
\label{app:reproducibility}
The accompanying archive contains:
\begin{description}[style=nextline]
\item[\RepoPath{rvachev\_lagrange\_loop.tex}]
The complete report source.
\item[\RepoPath{rvachev\_lagrange\_loop.pdf}]
The compiled report.
\item[\RepoPath{numerical\_experiments.py}]
Exact symbolic coefficient generation, FFT tabulation of \(u\), synthesis/right-inverse/projector checks, interpolation experiments, and figure/table generation.
\item[\RepoPath{prepare\_corpus.py}]
Recursive indexer for every TeX file in the repository corpus.
\item[\RepoPath{corpus\_manifest.txt} and \RepoPath{corpus\_manifest.csv}]
Human- and machine-readable source manifests.
\item[\RepoPath{relevant\_source\_excerpts.txt} and \RepoPath{lagrange\_source\_map.txt}]
Line-numbered contexts and a ranked map of files coupling Lagrange/interpolation terminology to the Rvachev function.
\item[\RepoPath{novelty\_overlap\_audit.txt}]
A conservative lexical overlap audit for the report's operator-level deductions and conjectural language.
\item[\RepoPath{*.csv}, \RepoPath{experiment\_summary.txt}, and \RepoPath{interpolation\_error.*}]
Raw numerical outputs and generated figure files.
\end{description}

To regenerate the numerical material and PDF from the report directory:
\begin{lstlisting}[language=bash]
python numerical_experiments.py
pdflatex -interaction=nonstopmode -halt-on-error rvachev_lagrange_loop.tex
pdflatex -interaction=nonstopmode -halt-on-error rvachev_lagrange_loop.tex
\end{lstlisting}

\backmatter
\begin{thebibliography}{99}

\bibitem{Fabius1966}
J.~Fabius,
\newblock A probabilistic example of a nowhere analytic \(C^\infty\)-function,
\newblock \emph{Zeitschrift f\"ur Wahrscheinlichkeitstheorie und Verwandte Gebiete} \textbf{5} (1966), 173--174.

\bibitem{StrangFix}
G.~Strang and G.~Fix,
\newblock A Fourier analysis of the finite element variational method,
\newblock in \emph{Constructive Aspects of Functional Analysis}, C.I.M.E., 1971, Edizioni Cremonese, Rome, 1973, 793--840.

\bibitem{deBoorRon1992}
C.~de Boor and A.~Ron,
\newblock Fourier analysis of the approximation power of principal shift-invariant spaces,
\newblock \emph{Constructive Approximation} \textbf{8} (1992), 427--462.

\bibitem{deBoorDeVoreRon1994}
C.~de Boor, R.~A.~DeVore, and A.~Ron,
\newblock The structure of finitely generated shift-invariant spaces in \(L_2(\mathbb R^d)\),
\newblock \emph{Journal of Functional Analysis} \textbf{119} (1994), 37--78.

\bibitem{BerrutTrefethen2004}
J.-P.~Berrut and L.~N.~Trefethen,
\newblock Barycentric Lagrange interpolation,
\newblock \emph{SIAM Review} \textbf{46} (2004), 501--517.

\bibitem{TrefethenATAP}
L.~N.~Trefethen,
\newblock \emph{Approximation Theory and Approximation Practice},
\newblock SIAM, Philadelphia, extended edition, 2019.

\bibitem{AriasDeReyna}
J.~Arias de Reyna,
\newblock An infinitely differentiable function with compact support: definition and properties,
\newblock online preprint and exposition on the Fabius--Rvachev function.

\bibitem{ProveItCorpus}
V.~Reshetnikov,
\newblock \emph{ProveIt: Analysis/FabiusFunction/docs}, recursive TeX corpus,
\newblock GitHub repository snapshot \RepoCommit, accessed 29 August 2026.

\bibitem{UpPolynomialSynthesis}
V.~Reshetnikov,
\newblock \emph{Up Polynomial Synthesis},
\newblock \RepoPath{Analysis/FabiusFunction/docs/Up\_Polynomial\_Synthesis.tex}, in \cite{ProveItCorpus}.

\end{thebibliography}

\end{document}
